\section{What the two sources contribute}\label{sec:sources}

\subsection{The upstream theorem and its transferable mechanism}

We inspect \texttt{lsernau/grh-lean} at commit
\texttt{72fead102f663d5c7a407964a13dc0cf9d977207} \cite{Sernau}.
Its exported statements concern RH and GRH for \emph{Dirichlet}
$L$-functions, including imprimitive characters. This scope does not
include every automorphic or Artin $L$-function. The repository reports
machine checks with the usual axioms \texttt{propext},
\texttt{Classical.choice}, and \texttt{Quot.sound}, using Lean 4.35.0 and
Mathlib commit \texttt{56a005855ea17e0b0958f73ec801ae140b0a1679}.
We inspected statements and selected proof components; we did not reproduce
the complete upstream build.

The arithmetic argument compares upper and lower bounds on the same
twisted lattice covolume. Divisibility produces short images, a second
moment produces sufficiently many independent images, and a determinant
argument controls their total size. Cancellation of the common main terms
yields a one-sided smoothed prime-sum bound, to which an analytic
zero-detection criterion applies. The component used here is the
finite-field estimate inside the second-moment argument.

In the compact source,
\texttt{grh-compact/src/New/Heights/OneDeterminant.lean},
the theorem \texttt{card\_mul\_sum\_det\_le} controls singular compressions
of a finite-group representation $\pi$:
\begin{equation}\label{eq:source}
 q\sum_{\det(A_1^{\tr}\pi(g)A_0)=0}\Theta(g)
 \le 2\sum_g\Theta(g)+q^{r+n/2}Y.
\end{equation}
Here the ambient dimension is $n$, the compression is $r$ by $r$,
$\Theta$ is nonnegative and of positive type, and $Y$ bounds the relevant
weighted stabilizer masses. Its proof uses additive-character
orthogonality, a tensor-product construction, and Cauchy--Schwarz.
It enters the Poincare-series second moment through
\texttt{grh-compact/src/New/Filling/Moments.lean}.
The readable counterpart is
\texttt{grh/src/New/Spanning/SecondMoment/PositiveType.lean}.

Our theorem determines the full reference distribution for the finite
symplectic group with uniform measure. It does not automatically replace
\eqref{eq:source} for arbitrary nonuniform positive-type weights.
Thus the connection is a concrete extension of a finite-field ingredient,
not a new consequence of the zero-location statement itself.

\subsection{The article's explicit gap}

The ProveIt \Article{} is pinned at commit
\texttt{fc9889c80312959aee9893ab21bfc9e741719e88} \cite{Article}.
Its README describes a merge of 103 reports into 621 pages and 760 stated
results, and says the merged mathematics was not independently recomputed.
We inspected the theorem inventory and the relevant proofs, not all 760
proofs. The archive records the section-level source hashes.

The section \texttt{spn.tex}, \emph{Spanning, frames, auxiliary characters
and matrix statistics}, proves a general compression bound, an exact
limiting law for an isotropic family, and an exact likelihood for an
arbitrary source with Lagrangian target. Its question
\texttt{spn:q:symplectic} asks for the law, or a limiting discrepancy, when
both compressed subspaces are arbitrary.
Theorem \ref{thm:law} answers the finite-law question for all dimensions
and restricted forms. Theorem \ref{thm:boundary} gives the limit in a
natural boundary regime with nondegenerate target.

\subsection{Contribution and prior art}

The proof uses classical finite geometry. Witt extension and symplectic
basis counting are standard; see Sprehn--Wahl \cite{SprehnWahl}.
De Boeck--Van de Voorde \cite{DeBoeck} give exact counts of nonsingular
subspaces disjoint from a fixed subspace and generating a nonsingular
span. Those incidence counts overlap special aggregations of our law.
Ishizuka \cite{Ishizuka} uses isometry-group orbits for uniform sampling
with prescribed hull dimension. These ingredients and symplectic incidence
counting are not claimed to be new.

The specific extension is a formula for \emph{each prescribed cross-pairing
matrix}, with arbitrary degeneracies on both sides, its three-invariant
support criterion, and the complement-Gram boundary law with quantitative
rank blindness. The novelty comparison is to the pinned \Article{}.
A focused primary-source search did not identify these statements in
this form; it does not establish worldwide priority.
